% from ega3-1-fr.tex
\oldpage[III]{153}

de $\mathfrak X$~; en d'autres termes, $\mathcal F$ est isomorphe au quotient d'un
$\mathcal O_{\mathfrak X}$-Module de la forme
$(\mathcal O_{\mathfrak X}(-n))^k$.
\end{corollary}

\begin{proof}
Comme $X_0$ est noethérien, il résulte de l'hypothèse sur $\mathcal L_0$ et de
(II, 4.5.5) qu'il existe $n_0$ tel que, pour $n\geq n_0$, $\mathcal F_0(n)$ soit
engendré par ses sections au-dessus de $\mathfrak X$~; par ailleurs, on peut supposer
$n_0$ pris assez grand pour que l'homomorphisme
\[
\Gamma(\mathfrak X,\mathcal F(n))\longrightarrow
\Gamma(\mathfrak X,\mathcal F_0(n))
\]
soit surjectif pour $n\geq n_0$ (\hyperref[III.5.2.3-fr]{5.2.3}). Il existe donc
un nombre fini de sections $s_i\in\Gamma(\mathfrak X,\mathcal F(n))$ dont les
images dans $\Gamma(\mathfrak X,\mathcal F_0(n))$ engendrent $\mathcal F_0(n)$
($0_{\mathrm I}$, 5.2.3). Comme $\mathfrak J$ est contenu dans l'idéal maximal de
l'anneau local en tout point de $\mathfrak X$, il résulte du lemme de Nakayama,
appliqué à ces anneaux locaux, que les $s_i$ engendrent $\mathcal F(n)$
($0_{\mathrm I}$, 5.1.1).
\end{proof}

\begin{env}[5.2.5]
\label{III.5.2.5-fr}
\emph{Démonstration du théorème d'existence : cas projectif.}
Les notations étant celles de (\hyperref[III.5.1.4-fr]{5.1.4}), supposons $f$
projectif, de sorte qu'il existe un $\mathcal O_X$-Module ample $\mathcal L$
(I, 5.5.4). Par définition, $X_n=\widehat X\times_{\widehat Y}Y_n$ est égal au
sous-préschéma fermé $X\times_Y Y_n=f^{-1}(Y_n)$ de $X$~; si
$\mathcal L'=\mathcal L\otimes_{\mathcal O_X}\mathcal O_{\widehat X}$ est le
complété de $\mathcal L$, on a donc
$\mathcal L'_0=\mathcal L/\mathfrak J\mathcal L$, considéré comme
$\mathcal O_{X_0}$-Module~; on sait que $\mathcal L'_0$ est ample
(II, 4.6.13, (i bis)). On peut donc appliquer à $\mathcal L'$ et à tout
$\mathcal O_{\widehat X}$-Module cohérent $\mathcal F$ le cor.
(\hyperref[III.5.2.4-fr]{5.2.4})~; on voit donc que $\mathcal F$ est isomorphe à
un quotient de
\[
\mathcal G=(\mathcal L'^{\otimes(-n)})^k
\]
pour des entiers $n>0$ et $k>0$ convenables. Or, il est clair que $\mathcal G$ est
le complété de $(\mathcal L^{\otimes(-n)})^k$ (I, 10.8.10), donc est
algébrisable. Considérons ensuite l'homomorphisme canonique surjectif
$u:\mathcal G\to\mathcal F$, et soit $\mathcal H=\operatorname{Ker}(u)$, qui est
un $\mathcal O_{\widehat X}$-Module cohérent ($0_{\mathrm I}$, 5.3.4). On voit de
même qu'il existe un $\mathcal O_{\widehat X}$-Module algébrisable $\mathcal K$ et
un homomorphisme $v:\mathcal K\to\mathcal G$ tel que
$\mathcal H=\operatorname{Im}(v)$. On a alors
$\mathcal F=\operatorname{Coker}(v)$, et $\mathcal F$ est algébrisable en vertu de
(\hyperref[III.5.2.2-fr]{5.2.2}).
\end{env}

\begin{env}[5.2.6]
\label{III.5.2.6-fr}
\emph{Démonstration du théorème d'existence; cas quasi-projectif.}
Les notations étant toujours celles de (\hyperref[III.5.1.4-fr]{5.1.4}), supposons
maintenant que $f$ soit quasi-projectif. Il existe alors un morphisme projectif
$g:Z\to Y$ tel que $X$ s'identifie au $Y$-préschéma induit sur un ouvert de $Z$
(II, 5.3.2)~; si on pose $Z'=g^{-1}(Y')$, on a $X'=X\cap Z'$. Par suite, le
complété $\widehat X=X_{/X'}$ s'identifie au préschéma formel induit par le complété
$\widehat Z=Z_{/Z'}$ sur l'ouvert $X\cap Z'$ de $\widehat Z$ (I, 10.8.5). Soit
$\mathcal F'$ un $\mathcal O_{\widehat X}$-Module cohérent, dont le support $T'$ est
propre sur $\widehat Y$~; cela signifie par définition qu'il existe un sous-préschéma
fermé de $X'$, ayant $T'\subset X'$ comme espace sous-jacent, tel que la restriction
$T'\to Y'$ de $f$ soit propre~; on en conclut que $T'$ est propre sur $Y$, donc
fermé dans $Z'$ (II, 5.4.10). Il en résulte que $\mathcal F'$ est le faisceau induit
sur $\widehat X$ par le $\mathcal O_{\widehat Z}$-Module $\mathcal G'$ obtenu par
recollement de $\mathcal F'$ (défini sur l'ouvert $\widehat X$ de $\widehat Z$) et
du faisceau $0$ sur l'ouvert $\widehat Z-T'$ de $\widehat Z$, ces deux faisceaux
coïncidant dans l'ouvert intersection $\widehat X-T'$. Il est clair que
$\mathcal G'$ est cohérent~; en vertu de
(\hyperref[III.5.2.5-fr]{5.2.5}), il existe un $\mathcal O_Z$-Module cohérent
$\mathcal G$ tel que $\mathcal G'=\widehat{\mathcal G}$~; soit $T$ le support de
$\mathcal G$, de sorte que $T'=T\cap Z'$ (I, 10.8.12). Si $h$ est la restriction
de $g$ au sous-préschéma fermé réduit de $Z$ ayant $T$ comme espace sous-jacent, on
a donc $T'=h^{-1}(Y')=T\cap g^{-1}(Y')$, et par suite $X\cap T$ est un ouvert de
$T$ contenant $T'$.

